% Préambule propre à la documentation SCONTO-SVU. Compilation avec XeLaTeX.
\usepackage[a4paper,margin=2.5cm]{geometry}
\usepackage{fontspec}
\usepackage[french]{babel}
\usepackage{microtype}
\usepackage{setspace}
\usepackage{titlesec}
\usepackage{graphicx}
\usepackage{float}
\usepackage{caption}
\usepackage{booktabs}
\usepackage{longtable}
\usepackage{tabularx}
\usepackage{array}
\usepackage{pdflscape}
\usepackage{xcolor}
\usepackage{hyperref}
\usepackage[nameinlink,noabbrev]{cleveref}
\crefname{chapter}{chapitre}{chapitres}
\Crefname{chapter}{Chapitre}{Chapitres}
\crefname{part}{partie}{parties}
\Crefname{part}{Partie}{Parties}
\crefname{section}{section}{sections}
\Crefname{section}{Section}{Sections}
\crefname{figure}{figure}{figures}
\Crefname{figure}{Figure}{Figures}
\crefname{table}{tableau}{tableaux}
\Crefname{table}{Tableau}{Tableaux}
\crefname{equation}{équation}{équations}
\Crefname{equation}{Équation}{Équations}
\providecommand{\crefrangeconjunction}{}
\providecommand{\crefpairconjunction}{}
\providecommand{\crefmiddleconjunction}{}
\providecommand{\creflastconjunction}{}
\renewcommand{\crefrangeconjunction}{ à }
\renewcommand{\crefpairconjunction}{ et }
\renewcommand{\crefmiddleconjunction}{, }
\renewcommand{\creflastconjunction}{ et }
\usepackage{fancyhdr}
\usepackage{lastpage}

\newcolumntype{L}[1]{>{\raggedright\arraybackslash}p{#1}}
\newcolumntype{Y}{>{\raggedright\arraybackslash}X}

\IfFontExistsTF{Times New Roman}{
  \setmainfont{Times New Roman}
}{
  \setmainfont{TeX Gyre Termes}
}

\IfFontExistsTF{Arial}{
  \newfontfamily\DocumentationSans{Arial}
}{
  \newfontfamily\DocumentationSans{TeX Gyre Heros}
}

\setstretch{1.12}
\setlength{\parindent}{0.6cm}
\setlength{\parskip}{0pt}

\titleformat{\chapter}[block]
  {\DocumentationSans\bfseries\fontsize{15}{18}\selectfont}
  {\thechapter.}{0.6em}{}[\vspace{0.25ex}\titlerule]
\titlespacing*{\chapter}{0pt}{-10pt}{16pt}

\titleformat{\section}
  {\DocumentationSans\bfseries\fontsize{12}{14}\selectfont}
  {\thesection}{0.6em}{}
\titleformat{\subsection}
  {\DocumentationSans\bfseries\fontsize{11}{13}\selectfont}
  {\thesubsection}{0.6em}{}

\captionsetup{font=small,labelfont=bf,justification=centering}

\hypersetup{
  colorlinks=true,
  linkcolor=black,
  citecolor=black,
  urlcolor=black,
  pdftitle={Documentation du modèle SCONTO-SVU},
  pdfauthor={}
}

\pagestyle{fancy}
\fancyhf{}
\fancyhead[L]{\small Documentation SCONTO-SVU}
\fancyhead[R]{\small Fonctionnement et utilisation}
\fancyfoot[C]{\small Page \thepage\ sur \pageref{LastPage}}
\renewcommand{\headrulewidth}{0.3pt}
\renewcommand{\footrulewidth}{0pt}

\definecolor{DocBlue}{HTML}{335C81}
\definecolor{DocBluePale}{HTML}{EDF3F8}
\definecolor{DocAmber}{HTML}{8A5A00}
\definecolor{DocAmberPale}{HTML}{FFF7E6}
\definecolor{DocGreen}{HTML}{356B4C}
\definecolor{DocGreenPale}{HTML}{EEF7F1}
\definecolor{DocGray}{HTML}{555555}
\definecolor{DocGrayPale}{HTML}{F3F3F3}

\newcommand{\DocumentationBox}[4]{%
  \par\smallskip\noindent
  \fcolorbox{#1}{#2}{%
    \begin{minipage}{0.92\linewidth}
      \small\textbf{#3}\quad #4
    \end{minipage}}%
  \par\smallskip
}

\newcommand{\BonASavoir}[1]{%
  \DocumentationBox{DocBlue}{DocBluePale}{Bon à savoir.}{#1}}
\newcommand{\AttentionDoc}[1]{%
  \DocumentationBox{DocAmber}{DocAmberPale}{Attention.}{#1}}
\newcommand{\ExempleDoc}[1]{%
  \DocumentationBox{DocGreen}{DocGreenPale}{Exemple.}{#1}}
\newcommand{\EnPratique}[1]{%
  \DocumentationBox{DocGray}{DocGrayPale}{En pratique.}{#1}}

\newif\ifdocumentationdraft
\documentationdrafttrue
\newcommand{\DocumentationNote}[1]{%
  \ifdocumentationdraft
    \par\smallskip\noindent
    \fcolorbox{DocGray}{DocGrayPale}{%
      \begin{minipage}{0.92\linewidth}
        \small\textit{Note interne de fondation : #1}
      \end{minipage}}
    \par\smallskip
  \fi
}

\newcommand{\FigureOrPlaceholder}[4][0.92\textwidth]{%
  \begin{figure}[H]
    \centering
    \IfFileExists{#2}{%
      \includegraphics[width=#1]{#2}%
    }{%
      \fbox{\parbox[c][4.5cm][c]{0.86\textwidth}{\centering\small
        Illustration planifiée : \texttt{\detokenize{#2}}}}%
    }%
    \caption{#3}
    \label{#4}
  \end{figure}%
}

\newcommand{\TablePlaceholder}[2]{%
  \begin{table}[H]
    \centering
    \begin{tabularx}{0.88\textwidth}{lY}
      \toprule
      Élément & Contenu attendu \\
      \midrule
      Fondation & Tableau à renseigner pendant la rédaction du chapitre. \\
      \bottomrule
    \end{tabularx}
    \caption{#1}
    \label{#2}
  \end{table}%
}
